\documentclass[12pt]{article}
\usepackage[utf8]{inputenc}
\usepackage[T1]{fontenc}
\usepackage{hyperref}
\title{Punan Merap Language Resources for K-12 Curriculum Development}
\author{Compiled by Claude Sonnet 5.5 (claude-sonnet-5-5) for Andr\'as Kornai}
\date{}
\begin{document}
\maketitle

\section{Introduction}
Punan Merap (Merap, ISO 639-3: puc; Glottocode puna1274) is an Austronesian language of Malinau Regency, North Kalimantan, Indonesia, reported in the village of Langap. Estimates are about 200 speakers (1981) to 300 (Joshua Project, citing Ethnologue). A journal study argues that ``Punan-Merap'' names a community of Punan living along the Merap River rather than a Punan-family language. It must not be confused with Punan Merah (puf), a separate language. No teaching materials were found.

\section{Primary Educational Resources}
None found.

\section{Linguistic and Reference Materials}
\begin{itemize}
\item Merap language (overview, phonology and classification): \url{https://en.wikipedia.org/wiki/Merap_language}
\item Endangered Languages Project entry for Punan Merap: \url{https://www.endangeredlanguages.com/lang/5298?hl=en}
\item ``A Holistic Study of Merap Language Mapping in Relation to the Term Punan-Merap'' (journal article): \url{https://jltr.academypublication.com/index.php/jltr/article/view/10062}
\item Punan Merah language (a different language, listed to avoid confusion): \url{https://en.wikipedia.org/wiki/Punan_Merah_language}
\end{itemize}

\section{Religious and Community Resources}
\begin{itemize}
\item Joshua Project profile of the Merap Punan people: \url{https://joshuaproject.net/people_groups/14480/ID}
\end{itemize}

\section{Digital and Audio Resources}
None found.

\section{Curriculum Development Context}
Only descriptive and demographic sources exist. A heritage or revitalization curriculum would have to start with documentation work with the Langap community, drawing on the phonological description cited on Wikipedia and the journal study above.
\end{document}